\chapter{Yoga: Meaning, Aim and History}
\label{ch:meaning-history}

\emph{What does the word yoga mean, what is yoga for, and how did it grow into the practice taught today?}
This chapter opens the theory. It explains the root of the word and the main classical definitions, the aim and objectives of yoga, the misconceptions a teacher meets in every class, and the broad periods through which yoga developed, with the texts and teachers a candidate should be able to name.

\section{The word and its senses}

The word \emph{yoga} comes from the Sanskrit root \emph{yuj}. The grammatical tradition credited to Panini gives the root three senses: to join or yoke (union), concentration (\emph{samādhi}), and control or restraint. These three senses explain why teachers define yoga in different ways, and a student who knows them is not confused by the variety. Patanjali uses the word mainly in the sense of concentration; devotional and Vedantic writers use the sense of union.

In the reading of union, yoga is the joining of the individual self (\emph{jīvātmā}) with the universal self (\emph{paramātmā}), or the merging of the apparent self in the true, limitless self. In a simpler reading, it is the harmony of body, mind and spirit, reached through mastery of the body, the senses and the mind.

The word is used in two ways, and a good answer says which is meant. It names the goal, the state reached (\emph{sādhya}), and it names the means, the practice that leads there (\emph{sādhana}).

\section{Classical definitions}

Give each definition with its source, and in plain words.
\begin{itemize}
\item \textbf{Yoga Sutra} (first chapter, second sutra): \emph{yogaś citta-vṛtti-nirodhaḥ}, yoga is the stilling of the movements of the mind.
\item \textbf{Bhagavad Gita} (second chapter): \emph{samatvaṃ yoga ucyate}, evenness of mind in success and failure is called yoga (verse 48); and \emph{yogaḥ karmasu kauśalam}, yoga is skill in action (verse 50). In the sixth chapter (verse 17) the Gita adds that yoga destroys sorrow for one who is moderate in food, recreation, effort, sleep and waking.
\item \textbf{Yoga Vasishtha}: yoga is the skilful means of calming the mind.
\item \textbf{Katha Upanishad}: yoga is the firm holding back of the senses.
\end{itemize}

\section{Aim and objectives}

The guidebook gives the ultimate aim of yoga as freedom of the individual from the sufferings of life, with a calm mind, will power and a balanced personality. In the traditional view the highest aim is liberation (\emph{kaivalya}, \emph{mukti}, \emph{mokṣa}); the Common Yoga Protocol joins this to holistic health, happiness and harmony. In public programmes, the stated aim is health and well-being.

For a class, a teacher can state the objectives simply: a healthy body, a calm and focused mind, emotional balance, better conduct towards others, and awareness of oneself. These are teaching summaries, not a Board list.

The guidebook ties a healthy and happy life to five supports: food (\emph{āhāra}), recreation and daily routine (\emph{vihāra}), conduct (\emph{ācāra}), thought (\emph{vicāra}) and behaviour towards others (\emph{vyavahāra}).

\section{Common misconceptions}

A teacher meets the same mistaken ideas again and again. Correct each with a short example from the practices.
\begin{itemize}
\item \emph{Yoga is only physical exercise, or only asanas.} It is holistic: it includes breath, attention, conduct and meditation.
\item \emph{You must be flexible first.} Practice builds flexibility, and breath and attention practices need none.
\item \emph{Yoga belongs to one religion.} It belongs to no single religion or community.
\item \emph{Yoga is for ascetics.} It is taught for everyday life.
\item \emph{Yoga replaces medical treatment.} It supports health and well-being; it does not replace treatment.
\item \emph{Yoga gives quick results.} Results come with regular practice over time.
\item \emph{Powers, magic and display are yoga.} The guidebook says that the pursuit of supernatural powers, magic shows and spectacles such as being buried alive are not yoga.
\item \emph{Yoga must be avoided in pregnancy.} The guidebook says doctors regard gentle yoga as helpful; a pregnant student should still ask a yoga expert first.
\end{itemize}

\section{Origin and periods}

Yoga is very ancient. The guidebook's foreword places its beginning almost 5,000 years ago and names the Rigveda as its first written reference. In yogic lore, Shiva is the first yogi (Adiyogi) and the first teacher (Adi Guru), who passed the teaching to the seven sages (\emph{saptarṣi}). Seals and figures from the Indus Saraswati Valley civilisation, dated to about 2700 BC in both the guidebook and the protocol, are cited as signs of early yoga practice; the guidebook names the Pashupati figure among them.

The guidebook groups the history into four broad periods:
\begin{itemize}
\item \textbf{Pre-Vedic}: traces in the Indus Valley civilisation, such as seals with seated figures.
\item \textbf{Vedic}: the Vedas and the Upanishads, in which practices of breath, austerity and meditation are described, and which turn the outer sacrifice into an inner search for self-knowledge.
\item \textbf{Classical}: the Yoga Sutra of Patanjali, which sets out the eightfold path in a short, systematic form; the epics and the Gita belong to the same broad stage.
\item \textbf{Post-classical}: Tantra and the hatha yoga texts, with their stress on postures, cleansing, breath control and locks.
\end{itemize}
Keep dates out of an answer unless a source gives them; name the periods and the texts.

\ycbfig{ycb1-fig-history}

\section{How yoga developed}

Yoga grew from a spiritual discipline found in several traditions, through its codification in the Yoga Sutra, to the body-centred hatha yoga, and then to the health-centred forms taught worldwide today. The texts carry this story, each adding something:
\begin{itemize}
\item the Vedas give the early ritual and philosophical base;
\item the Upanishads give the concepts and the stress on self-realisation;
\item the Yoga Sutra gives the eightfold path as a system;
\item the Bhagavad Gita joins yoga to life and duty, through karma, bhakti and jnana yoga;
\item the Hatha Pradipika gives the practical methods of hatha yoga.
\end{itemize}
Buddhist and Jain teaching, the epics and the philosophical systems also carry yoga ideas.

Modern times added many named styles. The guidebook mentions Power, Ashtanga and Vinyasa classes, and the protocol lists Kundalini yoga among the traditional schools. The loosening practices of sukshma vyayama, according to the guidebook, were developed by Maharshi Karthikeyaji Maharaj of the Himalaya and spread in the modern world by his student Swami Dhirendra Brahmachari.

\subsection*{Names to know}
\begin{itemize}
\item Patanjali, who codified classical yoga, and Vyasa, an early commentator on the Yoga Sutra.
\item Gorakshanath, of the hatha yoga tradition, and Svatmarama, author of the Hatha Pradipika.
\item Modern teachers who spread yoga in India and abroad, among them Swami Vivekananda (whose Chicago address of 1893 the guidebook takes as the start of the modern spread), T. Krishnamacharya, Swami Sivananda, Swami Kuvalayananda, B.K.S. Iyengar and Swami Satyananda.
\end{itemize}

\ycbfig{ycb1-fig-teachers}

\section{The philosophical schools}

Indian philosophy is divided into schools that accept the authority of the Vedas (\emph{āstika}) and schools that do not (\emph{nāstika}). The guidebook lists them with their founders:
\begin{itemize}
\item \textbf{Astika}: Nyaya (Gautama), Vaisheshika (Kanada), Samkhya (Kapila), Yoga (Patanjali), Purva Mimamsa (Jaimini) and Vedanta (Badarayana).
\item \textbf{Nastika}: Jain (Mahavira, the last of the twenty-four tirthankaras), Buddhist (the Buddha) and Charvaka (traditionally Brihaspati).
\end{itemize}

Yoga draws its theory from Samkhya. Samkhya teaches two eternal principles: \emph{puruṣa}, pure consciousness, and \emph{prakṛti}, matter, which is made of the three gunas (sattva, rajas and tamas). Liberation, \emph{kaivalya}, comes from discerning purusha from prakriti. It ends suffering of three kinds: from within oneself (\emph{ādhyātmika}), from other beings (\emph{ādhibhautika}) and from cosmic forces (\emph{ādhidaivika}). The gunas return in the chapter on the principles of yoga (see \cref{ch:principles-schools}).

\begin{recall}
\begin{itemize}
\item Root \emph{yuj}: union, concentration, control.
\item Patanjali: stilling of the movements of the mind. Gita: evenness of mind; skill in action.
\item Yoga is both the goal (sadhya) and the means (sadhana).
\item Periods: pre-Vedic, Vedic, classical, post-classical.
\item Astika: Nyaya, Vaisheshika, Samkhya, Yoga, Mimamsa, Vedanta. Nastika: Jain, Buddhist, Charvaka.
\end{itemize}
\end{recall}

\subsection*{Worked examples}
\begin{itemize}
\item \textbf{A new student says, ``I can't do yoga, I can't touch my toes.''} Explain that flexibility comes from practice, not before it, and that breath and attention practices need none. Start with seated breathing.
\item \textbf{Model answer: define yoga in three ways.} Patanjali: yoga is the stilling of the movements of the mind. The Gita: yoga is evenness of mind in success and failure, and skill in action. The Yoga Vasishtha: yoga is a skilful means of calming the mind. The differences follow from the three senses of the root \emph{yuj}.
\item \textbf{Model answer: the history of yoga in five lines.} Ancient origins, with the Rigveda as the first reference and Indus Valley seals as early signs; the Vedas and Upanishads; Patanjali's Yoga Sutra as the classical system; the hatha texts of the post-classical period; and the modern spread from the late nineteenth century.
\end{itemize}

\subsection*{Check yourself}
\begin{enumerate}
\item Which root does the word yoga come from, and what are its three senses?
\item Give Patanjali's definition in your own words.
\item Which text defines yoga as skill in action?
\item Which text calls yoga a skilful means of calming the mind?
\item What is the difference between sadhya and sadhana?
\item Name one misconception about yoga and correct it.
\item Which text sets out the eightfold path?
\item Which period is marked by the hatha yoga texts?
\item Who is the first yogi and first teacher in yogic lore?
\item Name the astika schools.
\item Which two principles does Samkhya teach?
\end{enumerate}

\subsection*{Answers}
\begin{enumerate}
\item \emph{Yuj}: union, concentration (samadhi) and control.
\item Yoga is the stilling of the movements of the mind.
\item The Bhagavad Gita.
\item The Yoga Vasishtha.
\item Sadhya is the goal; sadhana is the practice; the word yoga is used for both.
\item For example: ``yoga needs flexibility''; it does not, because flexibility comes with practice and breath and attention practices need none.
\item The Yoga Sutra of Patanjali.
\item The post-classical period.
\item Shiva (Adiyogi, Adi Guru).
\item Nyaya, Vaisheshika, Samkhya, Yoga, Mimamsa and Vedanta.
\item Purusha (consciousness) and prakriti (matter).
\end{enumerate}

\paragraph{Sources.} The official guidebook for Level I (2016, older scheme), first chapter; the Common Yoga Protocol, Ministry of AYUSH (introduction and history); the Yoga Sutra and the Bhagavad Gita (definitions); the Katha Upanishad with Shankara's commentary; the Board syllabus for Level 1.
